\documentclass[10pt,a4paper]{amsart}
\usepackage{amsmath,amssymb,amsthm,graphicx}
\usepackage{verbatim}
\newenvironment{reference}{\comment}{\endcomment}
%\newenvironment{reference}{}{}
\newenvironment{slogan}{\comment}{\endcomment}
\newenvironment{history}{\comment}{\endcomment}

% For commutative diagrams we use Xy-pic
\usepackage[all]{xy}

% We use 2cell for 2-commutative diagrams.
\xyoption{2cell}
% Theorem environments.
%
\theoremstyle{plain}
\newtheorem{theorem}[subsection]{Theorem}
\newtheorem{proposition}[subsection]{Proposition}
\newtheorem{lemma}[subsection]{Lemma}

\theoremstyle{definition}
\newtheorem{definition}[subsection]{Definition}
\newtheorem{example}[subsection]{Example}
\newtheorem{exercise}[subsection]{Exercise}
\newtheorem{situation}[subsection]{Situation}

\theoremstyle{remark}
\newtheorem{remark}[subsection]{Remark}
\newtheorem{remarks}[subsection]{Remarks}

\numberwithin{equation}{subsection}

% Macros
%
\def\lim{\mathop{\mathrm{lim}}\nolimits}
\def\colim{\mathop{\mathrm{colim}}\nolimits}
\def\Spec{\mathop{\mathrm{Spec}}}
\def\Hom{\mathop{\mathrm{Hom}}\nolimits}
\def\Ext{\mathop{\mathrm{Ext}}\nolimits}
\def\SheafHom{\mathop{\mathcal{H}\!\mathit{om}}\nolimits}
\def\SheafExt{\mathop{\mathcal{E}\!\mathit{xt}}\nolimits}
\def\Sch{\mathit{Sch}}
\def\Mor{\mathop{\mathrm{Mor}}\nolimits}
\def\Ob{\mathop{\mathrm{Ob}}\nolimits}
\def\Sh{\mathop{\mathit{Sh}}\nolimits}
\def\NL{\mathop{N\!L}\nolimits}
\def\CH{\mathop{\mathrm{CH}}\nolimits}
\def\proetale{{pro\text{-}\acute{e}tale}}
\def\etale{{\acute{e}tale}}
\def\QCoh{\mathit{QCoh}}
\def\Ker{\mathop{\mathrm{Ker}}}
\def\Im{\mathop{\mathrm{Im}}}
\def\Coker{\mathop{\mathrm{Coker}}}
\def\Coim{\mathop{\mathrm{Coim}}}

% Boxtimes
%
\DeclareMathSymbol{\boxtimes}{\mathbin}{AMSa}{"02}

%
% Macros for moduli stacks/spaces
%
\def\QCohstack{\mathcal{QC}\!\mathit{oh}}
\def\Cohstack{\mathcal{C}\!\mathit{oh}}
\def\Spacesstack{\mathcal{S}\!\mathit{paces}}
\def\Quotfunctor{\mathrm{Quot}}
\def\Hilbfunctor{\mathrm{Hilb}}
\def\Curvesstack{\mathcal{C}\!\mathit{urves}}
\def\Polarizedstack{\mathcal{P}\!\mathit{olarized}}
\def\Complexesstack{\mathcal{C}\!\mathit{omplexes}}
% \Pic is the operator that assigns to X its picard group, usage \Pic(X)
% \Picardstack_{X/B} denotes the Picard stack of X over B
% \Picardfunctor_{X/B} denotes the Picard functor of X over B
\def\Pic{\mathop{\mathrm{Pic}}\nolimits}
\def\Picardstack{\mathcal{P}\!\mathit{ic}}
\def\Picardfunctor{\mathrm{Pic}}
\def\Deformationcategory{\mathcal{D}\!\mathit{ef}}
\usepackage{cleveref}
\crefname{originalreference}{original reference}{original references}
\Crefname{originalreference}{Original reference}{Original references}
\makeatletter
\newlabel{lemma-finite-type-red}{{0GC2}{source}{}{sourcepointer-1}{}}
\newlabel{lemma-finite-type-red@cref}{{[originalreference][0][]0GC2}{[1][0][]source}{}{}{}}
\makeatother
\title{Finite type after reduction implies topologically of finite type}
\author{The Stacks Project Authors}
\date{stacks-project:restricted.tex:local-snapshot-20261012; GNU FDL 1.2 or later}
\begin{document}
\maketitle
\noindent\textit{Editorial note.} All delivered mathematical text is the author's own, copied from the pinned Stacks Project chapter file restricted.tex; no statement or proof step was rewritten, reordered or removed.\par\medskip
\noindent\textit{Editorial note.} The counted proof is the author's complete argument: the proof builds the candidate polynomial map from reduced generators plus ideal generators, expands elements of powers of b into images plus deeper powers by induction, and concludes density from the definition topology of b\par\medskip
\noindent\textit{Editorial note.} The label lemma-Noetherian-finite-type-red names the counted statement, and the proof environment immediately following it in the source is delivered verbatim as its proof.\par\medskip
\noindent\textit{Editorial note.} The counted statement is carried by the author's own lemma environment in the chapter source.\par\medskip
\noindent\textit{Allowed standard prerequisites (editorial).}\begin{itemize}
\item the condition of lemma-finite-type-red of the same chapter: B/b is generated over A by finitely many classes, where b is the ideal of topologically nilpotent elements
\item b is an ideal of definition of the weakly admissible Noetherian ring B
\item the definition of topologically of finite type via density of the image of a polynomial map
\end{itemize}\medskip
\ifdefined\hypertarget\hypertarget{sourcepointer-1}{}\fi
\noindent\textit{Source reference.} Original reference 0GC2: sources/stacks-restricted/pinned-original-source.tex, source lines 2631--2631. Prerequisite pointer invoked at proof line(s) 2751; it names a standard fact of the same chapter that the counted argument uses, and every step of the counted proof built on it is retained verbatim in the delivered body.\par\medskip
\begin{lemma}
\label{lemma-Noetherian-finite-type-red}
Let $\varphi : A \to B$ be an arrow of $\textit{WAdm}^{Noeth}$
satisfying the condition defined in Lemma \ref{lemma-finite-type-red}.
Then $A \to B$ is topologically of finite type.
\end{lemma}

\begin{proof}
Let $\mathfrak b \subset B$ be the
ideal of topologically nilpotent elements. Choose $b_1, \ldots, b_r \in B$
which map to generators of $B/\mathfrak b$ over $A$.
Choose generators $b_{r + 1}, \ldots, b_s$ of the ideal
$\mathfrak b$. We claim that the image of
$$
\varphi : A[x_1, \ldots, x_s] \longrightarrow B, \quad
x_i \longmapsto b_i
$$
has dense image. Namely, if $b \in \mathfrak b^n$ for some $n \geq 0$,
then we can write
$b = \sum b_E b_{r + 1}^{e_{r + 1}} \ldots b_s^{e_s}$ for multiindices
$E = (e_{r + 1}, \ldots, e_s)$ with $|E| = \sum e_i = n$ and $b_E \in B$.
Next, we can write $b_E = f_E(b_1, \ldots, b_r) + b'_E$
with $b'_E \in \mathfrak b$ and $f_E \in A[x_1, \ldots, x_r]$.
Combined we obtain $b \in \Im(\varphi) + \mathfrak b^{n + 1}$.
By induction we see that $B = \Im(\varphi) + \mathfrak b^n$ for all
$n \geq 0$ which mplies what we want as $\mathfrak b$ is an ideal
of definition of $B$.
\end{proof}

\end{document}
